\section{Experimental Setup}
\label{sec:experimental_setup}

\subsection{Base Models}
\label{sec:base_models}

We evaluate three open-weight model checkpoints. Two are Qwen2.5-7B \cite{qwen2024qwen25} and Llama3.1-8B \cite{grattafiori2024llama3}; the third is DeepSeek-R1-Distill-Llama-8B \cite{deepseek2025r1}. The DeepSeek checkpoint is distilled from Llama3.1-8B-Base and therefore shares the Llama lineage with Llama3.1-8B; the evaluated set should thus be interpreted as three checkpoints rather than three independent model families. For each checkpoint, \emph{Base} denotes the checkpoint before task-specific adaptation, \emph{DD-LoRA} denotes the checkpoint fine-tuned on the mixed corpus of general instructions and ordinary Dou Dizhu decision data, and \emph{RC-LoRA} denotes the counterpart trained on the constraint-enriched mixed corpus, in which 15.5\% of the Dou Dizhu examples contain explicit prohibitive constraints. Mechanistic experiments are restricted to Qwen2.5-7B, for which the complete single-head scan and Top-$k$ joint-ablation experiments were performed. Additional implementation and reproducibility details are provided in an online appendix.\footnote{\raggedright\samepage Online appendix: \url{https://github.com/farshelina/rule-constrained-task-transfer/blob/main/paper/appendix.pdf}.}

\subsection{Fine-Tuning Configuration}
\label{sec:finetuning_configuration}

All fine-tuned checkpoints use LoRA \cite{hu2022lora}. The pretrained parameters remain frozen, while trainable low-rank modules are inserted into the attention projections $q_{\mathrm{proj}}$, $k_{\mathrm{proj}}$, $v_{\mathrm{proj}}$, and $o_{\mathrm{proj}}$, and the feed-forward projections $\mathrm{gate}_{\mathrm{proj}}$, $\mathrm{up}_{\mathrm{proj}}$, and $\mathrm{down}_{\mathrm{proj}}$. LoRA is therefore applied to both attention and feed-forward projections, while the mechanistic analysis focuses specifically on attention heads. Table~\ref{tab:lora_config} summarizes the principal configuration used for the final rule-constrained runs and the mechanistic analyses.

\begin{table}[t]
\centering
\caption{Principal LoRA configuration used for the final rule-constrained runs.}
\label{tab:lora_config}
\begin{tabular}{ll}
\toprule
\rowcolor{NeutralCell}
Hyperparameter & Value \\
\midrule
LoRA rank & 16 \\
LoRA alpha & 32 \\
LoRA dropout & 0.05 \\
Learning rate & $1\times10^{-4}$ \\
Training epochs & 3 \\
Per-device batch size & 4 \\
Maximum sequence length & 2048 \\
Learning-rate scheduler & Cosine \\
Warmup ratio & 0.1 \\
Gradient checkpointing & Enabled \\
\bottomrule
\end{tabular}
\end{table}

\subsection{ACL Bench Evaluation}
\label{sec:acl_evaluation}

Cross-task transfer is evaluated on the 300-item ACL Bench introduced in Section~\ref{sec:acl_bench}. The benchmark contains 100 four-choice questions in each of three categories: Hierarchy Logic, Pattern Recognition, and Counter Logic.

\begin{figure}[ht]
\centering
\begin{tikzpicture}[
node distance=1cm,
box/.style={
draw,
rounded corners,
align=center,
minimum width=3cm,
minimum height=0.8cm
},
smallbox/.style={
draw,
rounded corners,
align=center,
minimum width=2.5cm,
minimum height=0.7cm
},
arrow/.style={
->,
thick
}
]


% Source domain

\node (source) [box]
{Dou Dizhu\\
Source Domain};


% Abstract operations

\node (hier) [smallbox, below left=1.2cm and 2.0cm of source]
{
Hierarchical\\
Relations
};

\node (pattern) [smallbox, below=1.2cm of source]
{
Pattern\\
Structures
};

\node (counter) [smallbox, below right=1.2cm and 2.0cm of source]
{
Conditional\\
Selection\\
Operations
};


\draw[arrow] (source) -- (hier);
\draw[arrow] (source) -- (pattern);
\draw[arrow] (source) -- (counter);



% ACL Bench

\node (acl) [box, below=2cm of pattern]
{
ACL Bench\\
Abstract Evaluation
};


\node (logic1) [smallbox, below left=1cm and 1.8cm of acl]
{
Hierarchy\\
Logic
};


\node (logic2) [smallbox, below=1cm of acl]
{
Pattern\\
Recognition
};


\node (logic3) [smallbox, below right=1cm and 1.8cm of acl]
{
Counter\\
Logic
};


\draw[arrow] (hier) -- (acl);
\draw[arrow] (pattern) -- (acl);
\draw[arrow] (counter) -- (acl);


\draw[arrow] (acl) -- (logic1);
\draw[arrow] (acl) -- (logic2);
\draw[arrow] (acl) -- (logic3);


\end{tikzpicture}
\caption{
Construction of ACL Bench by abstracting Dou Dizhu decision operations into transferable reasoning categories.
The benchmark removes domain-specific surface forms while preserving underlying decision operations.
}
\label{fig:aclbench}
\end{figure}

We report overall and category-level accuracy for Base, DD-LoRA, and RC-LoRA. This three-way comparison separates the effect associated with ordinary Dou Dizhu mixed-data adaptation from the additional change observed in the constraint-enriched condition. Category-level results are retained because fine-tuning effects differ across reasoning types even when overall accuracy improves.

\subsection{Constraint-Sensitive Probe Evaluation}
\label{sec:constraint_probe_setup}

Constraint-induced relative preference shifts are evaluated with 60 controlled probes covering six decision scenarios. Each probe has a clean condition and a constrained condition in which an otherwise preferred action is explicitly prohibited. Evaluation uses seed 42, 16-bit floating-point (FP16) inference, scaled dot-product attention (SDPA), and batch size 1. We report mean BCDM across all 60 probes and per-scenario scores, using the metric defined in Section~\ref{sec:constraint_metric}; the six scenarios are reported in Table~\ref{tab:constraint_scenarios}. These probes are used to measure relative preference shifts under explicit prohibition rather than as an independent benchmark or a direct measure of generation compliance.

\subsection{Single-Head Causal Analysis}
\label{sec:causal_setup}

The complete attention-head scan is performed on Qwen2.5-7B, which has 28 Transformer layers and 28 attention heads per layer, for a total of 784 heads. Each head is ablated independently and the model is reevaluated on the same 60 constraint-sensitive probes. The scan is run for both Base and RC-LoRA, yielding 1,568 single-head intervention evaluations. For each head $(l,h)$, we compute $CE_{\mathrm{Base}}(l,h)$, $CE_{\mathrm{RC\text{-}LoRA}}(l,h)$, and $\Delta CE(l,h)$ as defined in Section~\ref{sec:head_ablation}. The resulting $28\times28$ map is used to analyze causal changes and rank heads for joint ablation.

\subsection{Top-k Joint-Ablation Evaluation}
\label{sec:topk_setup}

Joint ablation is evaluated at $k\in\{5,10,27\}$. For each $k$, we compare three conditions: \emph{Top-$k$}, consisting of the heads with the largest positive $\Delta CE$ values; \emph{Global Random}, sampled from all 784 heads; and \emph{Matched Random}, sampled to preserve the coarse layer-bin distribution of the corresponding Top-$k$ set across Layers 0--6, 7--13, 14--20, and 21--27. All joint-ablation runs use the same 60 probes, seed 42, FP16 inference, SDPA attention, and batch size 1. For each random-control condition, five independently sampled sets are generated with seeds 0--4, and their degradation scores are averaged.

The two random controls separate generic sensitivity to multi-head removal from effects attributable to broad depth-region concentration. Because the Matched Random condition is bin-matched rather than exact-layer matched, it does not eliminate all possible sensitivity differences among individual layers. A larger degradation under Top-$k$ than under both controls therefore supports additional functional relevance of the selected head identities beyond intervention size and coarse depth distribution. No additive assumption is made for joint interventions.

\subsection{Implementation Environment}
\label{sec:implementation_environment}

Fine-tuning and mechanistic experiments are conducted on an NVIDIA RTX A4500 GPU with 20~GB of memory. Gradient checkpointing is enabled during LoRA fine-tuning. Mechanistic evaluation uses mixed-precision inference where applicable and batch size 1 to support consistent head-level intervention.

